\documentclass{article}

\usepackage{amsmath,amssymb}
\usepackage{xurl}
\usepackage[hidelinks]{hyperref}
\setlength{\emergencystretch}{2em}

\input{command}

\title{Formalization-Derived Corrections to Wolf's \emph{Quantum Channels \& Operations}}
\author{}
\date{2026-08-24}

\begin{document}
\sloppy
\maketitle
\gapnote{local-correction}{resolved}

\begin{abstract}
This note records only discrepancies in Michael M.~Wolf's \emph{Quantum
Channels \& Operations: Guided Tour} that were exposed by the present Lean
formalization or by a paper-gap analysis required to state a formal theorem
faithfully. It is not a general errata list for the lecture notes. Its main
body records only formalization-derived corrections; the final appendix records
three separately audited local PDF--transcription errors. Each item quotes the
relevant original statement, gives the locally correct replacement, and explains
why the correction is necessary.
\end{abstract}

\section{Scope}

The source PDFs are stored in \path{Notes/WolfNotePDF/}; the local
transcription is in \path{Notes/WolfNoteTexSource/}. The corresponding
formalization and paper-gap records are named explicitly below. A distinction
is essential:
\begin{itemize}
  \item a \emph{source error} is a printed mathematical assertion that is false
    as stated;
  \item a \emph{well-posedness correction} makes an implicit boundary convention
    explicit, without changing the substantive theorem;
  \item a \emph{formalization scope gap} is not an erratum: it records that the
    current Lean theorem is narrower than the source or that an auxiliary proof
    remains incomplete.
\end{itemize}

Only the first two kinds belong in the formalization-derived main body. The
appendix keeps transcription errors visibly separate from that record.

\section{The trace sign in the Lorentz-cone converse}
\label{sec:lorentz}

\paragraph{Formalization trigger.}
The formal proof of the converse trace inequality requires nonnegativity of the
trace. Its declaration is
\path{Matrix.IsHermitian.posSemidef_of_trace_re_nonneg_of_card_sub_one_mul_trace_sq_re_le}
in \doclink{QICLean/Analysis/MatrixTraceInequalities}. The discrepancy is
tracked in \path{docs/paper-gaps/wolf_ch3_lorentz_cone_trace_sign.tex}.

\paragraph{Printed source and its role.}
The proposition occurs at the end of the discussion of automorphisms of the
positive cone in Chapter~3.  Its first part gives the elementary necessary
inequality for a positive semidefinite Hermitian matrix; the second part is
intended to characterize the corresponding Lorentz cone.  Chapter~3,
Proposition~3.7 in the local PDF (book p.~66) states, for Hermitian $A$,
\begin{equation}\tag{printed converse}
  \operatorname{tr}(A)^2\geq(d-1)\operatorname{tr}(A^2)
  \quad\Longrightarrow\quad A\geq0.
\end{equation}

\paragraph{Correction.}
The valid converse is
\begin{equation}\tag{formalized converse}
  A=A^\dagger,\qquad
  \operatorname{tr}(A)\geq0,\qquad
  \operatorname{tr}(A)^2\geq(d-1)\operatorname{tr}(A^2)
  \quad\Longrightarrow\quad A\geq0.
\end{equation}
In the complex-matrix Lean statement the traces are expressed through real
parts; for Hermitian matrices these are the same real quantities.

\paragraph{Why the correction is necessary.}
Take $A=-I_d$. It is Hermitian and not positive semidefinite, but
\begin{equation}
  \operatorname{tr}(A)^2=d^2
  \geq d(d-1)=(d-1)\operatorname{tr}(A^2).
\end{equation}
Thus the squared inequality alone describes both orientations of the Lorentz
cone. The condition $\operatorname{tr}(A)\geq0$ selects the positive sheet.
It is also used by the proof: when a negative eigenvalue is present, the source
bounds $\operatorname{tr}(A)$ above by a nonnegative quantity and then squares
that bound. Squaring is justified only after fixing the sign of the trace.

\paragraph{Classification.}
This is a genuine source error. The unrestricted printed converse should not
be claimed as formalized; the formal theorem proves the corrected,
trace-nonnegative statement.

\section{The missing nonvanishing condition in Kramers' theorem~II}
\label{sec:kramers-ii}

\paragraph{Formalization trigger.}
The corrected theorem is formalized as
\leanid{Matrix.IsHermitian.two_le_finrank_eigenspace_of_intertwiner_mulVec_star_ne_zero}
(conditional form) and
\leanid{Matrix.IsHermitian.two_le_finrank_eigenspace_of_antisymmetric_isUnit}
(invertible intertwiner) in \doclink{QICLean/Algebra/KramersDegeneracy}; the
first Kramers theorem is recovered there as the unitary special case. The
source error is documented in
\path{docs/paper-gaps/wolf_ch3_kramers_theorem_ii.tex}.

\paragraph{Printed source and its role.}
After Kramers' theorem (a Hermitian $H$ commuting with an anti-unitary $T$,
$T^{2}=-\Id$, has all eigenvalues at least two-fold degenerate), Wolf prints
``Kramer's theorem II'': if $H$ is Hermitian and $HA=AH^{T}$ for some
anti-symmetric $A\neq 0$, then each eigenvalue of $H$ is at least two-fold
degenerate, ``as the same proof shows.''  The local source is
\path{Notes/WolfNoteTexSource/ch03_positive_not_completely.tex},
lines~527--530, with the matrix reduction at lines~500--501.

\paragraph{Correction.}
The proof transports an eigenvector $\ket\psi$ of eigenvalue $\lambda$ to the
partner $A\ket{\overline\psi}$, which is again a $\lambda$-eigenvector and is
orthogonal to $\ket\psi$; both steps survive for a general anti-symmetric
intertwiner.  The degeneracy follows only when the partner is nonzero, and
that is an extra condition: the correct local statement concludes
$\dim\ker(H-\lambda)\ge 2$ from $A\ket{\overline\psi}\neq 0$, and a global
conclusion requires $A$ injective.  The corrected theorem therefore assumes
$A$ invertible and concludes that every eigenvalue of $H$ is at least
two-fold degenerate.

\paragraph{Why the correction is necessary.}
On $\C^{3}$, $H=\operatorname{diag}(0,1,1)$ and the anti-symmetric unit of
the $1$-eigenspace satisfy $HA=AH^{T}=A$ with $A\neq 0$, but the eigenvalue
$0$ is simple: the partner $Ae_{1}$ of its eigenvector vanishes.

\paragraph{Classification.}
This is a genuine source error.  Theorem~I is unaffected: it is the case
$A=V^{\dagger}$ of the corrected theorem, where unitarity supplies exactly the
missing nonvanishing.

\section{The false converse for Pauli-block truncation}
\label{sec:pauli-block-truncation}

\paragraph{Formalization trigger.}
The forward preservation properties of Pauli-block truncation are formalized as
\leanid{Wolf.pauliBlockDiagonalTruncation_isPositiveMap} and
\leanid{Wolf.pauliBlockDiagonalTruncation_isCPMap} in
\doclink{QICLean/Channel/LorentzNormalForm/PauliBlockTruncation}. The source
error is documented in
\path{docs/paper-gaps/wolf_prop2_10_block_truncation_converse.tex}.

\paragraph{Printed source and its role.}
In Section~2.4, after the Bloch-ball representation in Eq.~(2.39), Wolf's
Proposition~2.10 considers the Pauli transfer matrix
\begin{equation}
  \widehat T=
  \begin{pmatrix}
    \widehat T_{00}&r^{\mathsf T}\\
    v&\Delta
  \end{pmatrix}
\end{equation}
of a Hermiticity-preserving map and prints the assertion that the block
truncation
\begin{equation}
  \widehat T'=\widehat T_{00}\oplus\Delta
\end{equation}
is positive, respectively completely positive, if and only if $\widehat T$ is.
The local source is
\path{Notes/WolfNoteTexSource/ch02_representations.tex}, lines~984--990.

\paragraph{Correction.}
With $\Theta(X)=\operatorname{tr}(X)\Id-X$ and
$D=\operatorname{diag}(1,-1,-1,-1)$, the truncation is
\begin{equation}
  T'=\frac{T+\Theta\circ T\circ\Theta}{2},
  \qquad
  \widehat T'=\frac{\widehat T+D\widehat T D}{2}
  =\widehat T_{00}\oplus\Delta.
\end{equation}
The valid statements are the forward implications
\begin{equation}
  T\text{ positive}\Longrightarrow T'\text{ positive},
  \qquad
  T\text{ completely positive}\Longrightarrow
  T'\text{ completely positive}.
\end{equation}
These are exactly the directions proved in the source at lines~992--998 and
formalized by the declarations above.

\paragraph{Why the correction is necessary.}
Let
\begin{equation}
  T(X)=\operatorname{tr}(X)\operatorname{diag}(3/2,-1/2).
\end{equation}
Then $T$ is Hermiticity preserving but not positive, since
$T(\Id/2)=\operatorname{diag}(3/2,-1/2)$. Its Pauli-block truncation is the
completely depolarizing channel $T'(X)=\operatorname{tr}(X)\Id/2$, which is
completely positive. Thus both printed converses fail.

\paragraph{Classification.}
This is a genuine source error. The proof of Proposition~2.10 supports the two
forward implications, not the printed equivalences.

\section{The inverse coefficient in the proof of Theorem~2.3}
\label{sec:ordered-cp-inverse}

\paragraph{Formalization trigger.}
The formal proof defines Wolf's normalized auxiliary matrix as
\leanid{stinespringW} and proves that the assignment from a supplied
Stinespring matrix is injective in \leanid{stinespringW_injective}, both in
\doclink{QICLean/Channel/OrderedCP}.

\paragraph{Printed source and its role.}
In the proof of Theorem~2.3, Wolf takes the normalized vector
$\ket{\Omega}=d'^{-1/2}\sum_a\ket{a,a}$ and defines
\begin{equation}
  W_i=(\one_{r_i}\otimes\bra{\Omega})
      (V_i\otimes\one_{d'}).
\end{equation}
The footnote asserting that $V_i\mapsto W_i$ is one-to-one prints the inverse
with the coefficient $d'^2$. The same coefficient appears in the local source
at \path{Notes/WolfNoteTexSource/ch02_representations.tex}, lines~392--393,
and in the PDF on book p.~39.

\paragraph{Correction.}
The coordinate relation is
\begin{equation}
  (W_i)_{j,(k,a)}=d'^{-1/2}(V_i)_{(a,j),k},
\end{equation}
so its coordinate inverse multiplies by $\sqrt{d'}$. Equivalently, with the
canonical tensor-factor identifications used in the source, the inverse is
\begin{equation}
  V_i=d'(W_i\otimes\one_{d'})(\one_d\otimes\ket{\Omega}).
\end{equation}

\paragraph{Why the correction is necessary.}
The second occurrence of the normalized vector contributes another factor
$d'^{-1/2}$. Hence the tensor expression without its leading coefficient is
$V_i/d'$, and the leading coefficient must be $d'$, not $d'^2$. The theorem
itself is unaffected: the formal injectivity proof uses the coordinate
relation and only cancels the nonzero factor $d'^{-1/2}$.

\paragraph{Classification.}
This is a normalization typo in a source footnote. It does not change the
statement or proof route of Theorem~2.3.

\section{The nonempty-family boundary in the Radon--Nikodym theorem}
\label{sec:radon-nikodym}

\paragraph{Formalization trigger.}
The source-facing Lean theorem
\leanid{IsKrausCP.radon_nikodym_of_stinespring} makes the finite index type
nonempty. It has Wolf's rectangular Heisenberg dimensions
$T_i,T:\mathcal M_{d'}\to\mathcal M_d$ and
$V:\mathbb C^d\to\mathbb C^{d'}\otimes\mathbb C^r$. The boundary is
documented in
\path{docs/paper-gaps/wolf_radon_nikodym_nonempty_family.tex}.

\paragraph{Printed source and its role.}
This theorem is the immediate instrument-valued corollary of the preceding
comparison theorem for ordered completely positive maps.  A supplied
Stinespring representation of the total map is held fixed, and the theorem
resolves each CP component by a positive effect on the same dilation space.
The exact local source is
\path{Notes/WolfNoteTexSource/ch02_representations.tex}:399--407. Chapter~2,
Theorem~2.4 (book p.~39) says:
\begin{quote}
``Let $\{T_i\}$ be a set of completely positive linear maps such that
$\sum_iT_i=T$ \ldots\ Then there is a set of positive operators $P_i\in M_r$
which sum up to $I_r$ such that
$T_i(A)=V^\dagger(A\otimes P_i)V$.''
\end{quote}

\paragraph{Formal statement.}
The formal theorem uses a \emph{nonempty finite family} of completely positive
linear maps
$T_i,T:\mathcal M_{d'}\to\mathcal M_d$ with
$\sum_{i\in\iota}T_i=T$, together with a supplied map
$V:\mathbb C^d\to\mathbb C^{d'}\otimes\mathbb C^r$ satisfying
$T(A)=V^\dagger(A\otimes I_r)V$. It concludes that there are positive
semidefinite operators $P_i\in\mathcal M_r$ such that
\begin{equation}
  \sum_{i\in\iota}P_i=I_r,
  \qquad
  T_i(A)=V^\dagger(A\otimes P_i)V.
\end{equation}
The declaration \leanid{IsCPMap.radon_nikodym_of_stinespring} is retained as
the square-algebra specialization.

\paragraph{Formal proof path.}
Wolf gives no separate proof beyond calling the result a simple corollary of
the ordered-map comparison theorem. The formal proof makes this corollary
explicit. It combines Kraus families for the components into an
outcome-labelled Stinespring representation $\widehat V$ of $T$, then applies
the preceding theorem to $T\leq T$ and the two representations
$\widehat V,V$. This gives a contraction, not in general an isometry,
\begin{equation}
  \widehat V=(I_{d'}\otimes C)V,
  \qquad C^\dagger C\leq I_r.
\end{equation}
If $C_i$ retains the rows labelled by $i$, the preliminary effects
$E_i=C_i^\dagger C_i$ are positive, represent $T_i$, and sum to
$C^\dagger C$. The residual
$R=I_r-C^\dagger C$ is positive and satisfies
$V^\dagger(A\otimes R)V=0$. A chosen outcome $i_0$ receives this residual,
so that $P_{i_0}=E_{i_0}+R$ and $P_i=E_i$ otherwise. This last step is needed
for an arbitrary nonminimal supplied dilation and is exactly where
nonemptiness is used. When $d'=0$, the formal proof separately assigns $I_r$
to a chosen outcome and zero to all others. The square-algebra declaration is
a direct specialization of the rectangular theorem.

\paragraph{Why the boundary is needed.}
If the index type is empty, then $\sum_iT_i=T$ forces $T=0$. Taking the zero
Stinespring matrix $V=0$ still satisfies the representation hypothesis, but the
printed conclusion would require
\begin{equation}
  0=\sum_iP_i=I_r,
\end{equation}
which is false when the dilation space has nonzero dimension. For the intended
notion of a quantum instrument, a nonempty finite outcome family is the
substantive case.

\paragraph{Classification.}
This is a local well-posedness correction, not an assertion that the
nonempty-case Radon--Nikodym theorem is false. The source's word ``set'' and
summation notation leave the boundary implicit; Lean must state it.

\section{Suppressed zero multiplicities in the two-pure-state spectrum}
\label{sec:two-pure-state-spectrum}

\paragraph{Formalization trigger.}
The Lean theorem
\leanid{Projectivization.charpoly_add_pureStateMatrix_of_two_le} keeps the full
characteristic polynomial of a sum of two normalized pure-state matrices. The
correction is documented in
\path{docs/paper-gaps/wolf_ch01_two_pure_state_zero_multiplicities.tex}.

\paragraph{Printed source and its role.}
In the spectrum-preserver argument following Wigner's theorem, Chapter~1 says
that the spectrum of two Hermitian rank-one projections $P,Q$ is
\begin{equation}
  1\pm\sqrt{\operatorname{tr}(PQ)}.
\end{equation}
These values are the two roots of the quadratic
$(X-1)^2-\operatorname{tr}(PQ)$.

\paragraph{Correction.}
For $d\geq2$, the multiplicity-aware statement is
\begin{equation}
  \chi_{P+Q}(X)
  =X^{d-2}\bigl((X-1)^2-\operatorname{tr}(PQ)\bigr).
\end{equation}
The factor $X^{d-2}$ supplies the guaranteed additional zero multiplicity.
If $P=Q$, then the quadratic is $X(X-2)$ and contributes one further zero, so
the total zero multiplicity is $d-1$. Dimension one has the separate formula
$\chi_{P+Q}(X)=X-2$, and dimension zero has no projective pure states.

\paragraph{Classification.}
This is a local correction to the displayed spectrum shorthand, not a claim
that Wolf's intended spectrum-preserver argument is false. The two displayed
values contain the transition-probability information used in that argument;
the characteristic polynomial records the complete multiplicity data.

\section{The rotation sign in the spinor exponential formula}
\label{sec:spinor-rotation-sign}

\paragraph{Formalization trigger.}
The rotation exponential is formalized as
\leanid{Wolf.spinorMatrix\_rotationExpSL2} in
\doclink{QICLean/Channel/LorentzNormalForm/SpinorExponential}.  The discrepancy
is documented in
\path{docs/paper-gaps/wolf_ch2_spinor_rotation_sign.tex}.

\paragraph{Printed source and its role.}
Wolf's Equation~(2.44), in
\path{Notes/WolfNoteTexSource/ch02_representations.tex}, lines~1070--1077,
defines
\begin{equation}
  R_i=\sum_{j,k=1}^3\varepsilon_{ijk}\ketbra{k}{j}
\end{equation}
and prints
\begin{equation}
  e^{-i\vec n\cdot\vec\sigma/2}\longmapsto e^{-\vec n\cdot\vec R}.
\end{equation}

\paragraph{Correction.}
With the standard Pauli matrices and column-vector convention used in the
notes, the corrected correspondence is
\begin{equation}
  e^{-i\vec n\cdot\vec\sigma/2}\longmapsto e^{+\vec n\cdot\vec R}.
\end{equation}
Equivalently, the printed Lorentz exponential corresponds to the spinor with
opposite sign in its exponent.

\paragraph{Why the correction is necessary.}
For \(\vec n=(0,0,\theta)\), direct Pauli conjugation gives
\begin{equation}
  e^{-i\theta\sigma_3/2}\sigma_1e^{i\theta\sigma_3/2}
   =\cos\theta\,\sigma_1+\sin\theta\,\sigma_2.
\end{equation}
The printed generator satisfies \(R_3\ket{1}=\ket{2}\), so this is the action
of \(e^{+\theta R_3}\), whereas \(e^{-\theta R_3}\) rotates toward the
negative second coordinate.

\paragraph{Classification.}
This is a genuine source sign error.  The formal theorem follows the actual
congruence action; the boost formula in the same equation is unchanged.

\section{The two boundary errors in Theorem~6.3}
\label{sec:theorem-6-3}

\paragraph{Formalization trigger.}
The declarations
\leanid{exists_wolfTheorem63_of_irreducible_positive} and
\leanid{exists_wolfTheorem63_of_irreducible_positive_of_ne_zero} formalize
Wolf Theorem~6.3 for arbitrary positive maps. The declarations
\leanid{lowerCollatzWielandtValue},
\leanid{upperCollatzWielandtValue}, and
\leanid{lowerCollatzWielandtGlobal\_eq\_upperCollatzWielandtGlobal}
formalize the corrected pointwise and global quantities. The two corrections
are documented in
\path{docs/paper-gaps/wolf_thm6_3_positive_map_cp_scope.tex}.

\paragraph{The global upper quantity.}
At local source lines~617--619, after defining the pointwise lower and upper
functionals $r(X)$ and $\widetilde r(X)$, the notes print
\begin{equation}\tag{printed}
  r=\sup_{X\geq0}r(X),\qquad
  \widetilde r=\sup_{X\geq0}\widetilde r(X),\qquad
  r\geq\widetilde r.
\end{equation}
The corrected Collatz--Wielandt pair is
\begin{equation}\tag{corrected}
  r=\sup_{X\geq0,\,X\neq0}r(X),\qquad
  \widetilde r=\inf_{X\geq0,\,X\neq0}\widetilde r(X).
\end{equation}
For an irreducible positive map, both extrema are attained at the same
positive-definite eigenvector and are equal. Thus the second ``sup'' and the
displayed reversed inequality are two local source errors. Equality is
not asserted pointwise for arbitrary $X$. In the extended-real formulation,
$r(0)=+\infty$ and $\widetilde r(0)=-\infty$; for every nonzero $X\geq0$ one
has $r(X)\leq\widetilde r(X)$, with equality at positive eigenvectors as used
at local source lines~649--651.

\paragraph{The one-dimensional zero map.}
The theorem assumes only that $T:\MN{D}\to\MN{D}$ is irreducible and positive,
but item~2 and the proof require a strictly positive Perron value. On
$\MN{1}$, the zero map is positive and irreducible because $0$ and $\Id$ are
the only orthogonal projections. Its spectral radius is zero. Hence the
assumption-free conclusion must allow $r=0$ while retaining the
positive-definite eigenvector $X>0$. Wolf's printed $r>0$ form becomes valid
after adding $T\neq0$. For $D>1$, irreducibility itself excludes the zero map.

\paragraph{Classification.}
These are two local false-source corrections exposed by the positive-map
formalization. They do not alter the proof for a nonzero irreducible map once
the second global extremum is read as an infimum.

\section{The missing order hypothesis in the unitary comparison}
\label{sec:unitary-comparison}

\paragraph{Formalization trigger.}
The correction and the exact source proof route are documented in
\path{docs/paper-gaps/wolf_ch5_unitary_comparison_missing_order.tex}.  The
corrected theorem is not yet formalized because the current dependency surface
does not expose the required cross-matrix Weyl monotonicity theorem.

\paragraph{Original statement and correction.}
Immediately after stating Weyl monotonicity for Hermitian matrices satisfying
$A\le B$, Wolf's Proposition~5.3 drops that hypothesis and claims that for
arbitrary Hermitian $A,B$ there is a unitary $U$ such that
\begin{align}
  f(A)\le U f(B)U^\dagger
\end{align}
for every non-decreasing scalar function $f$ on an interval containing the two
spectra.  In dimension one, $A=[1]$, $B=[0]$, and $f(x)=x$ contradict the
printed claim.  The corrected proposition inserts $A\le B$.

\paragraph{Source proof route.}
Under $A\le B$, Equation~(5.56) gives
$\lambda_j^\downarrow(A)\le\lambda_j^\downarrow(B)$ for every $j$ because
the smallest eigenvalue of $B-A$ is nonnegative.  Monotonicity of $f$ preserves
these scalar inequalities.  Choosing decreasingly ordered eigenvector
unitaries $V_A,V_B$ and setting $U=V_A V_B^\dagger$ then transports the
entrywise diagonal inequality to
$f(A)\le Uf(B)U^\dagger$.  This unitary depends only on the two ordered
eigenbases, so it works simultaneously for every admissible $f$, as required
by the quantifier order of the printed proposition.

\paragraph{Classification.}
This is a genuine source error: the printed theorem is false, while its
intended corrected statement follows by exactly the preceding argument once
$A\le B$ is retained.

\section{Formalization gaps that are not source corrections}

The following Wolf-specific paper-gap notes are intentionally \emph{not}
errata entries:
\begin{itemize}
  \item \path{wolf_prop28_square_case.tex}: the current normal-form theorem is
    the square-dimensional special case of Wolf's rectangular statement;
  \item \path{wolf_lemma_3_1_top_index_scope.tex},
    \path{wolf_prop_3_2_top_index_scope.tex}, and
    \path{wolf_t_eta_top_index_scope.tex}: endpoint Schmidt-rank cases that
    were formalization scope gaps and are now resolved;
  \item \path{wolf_reduction_criterion_schmidt_premise.tex}: the
    Schmidt-number premise was a formalization dependency and is now resolved;
  \item \path{wolf_ex3_1_choi_positivity_subcase_scope.tex}: the development
    proves selected parameter families, not the full source range;
  \item \path{wolf_ch5_operator_jensen_lieb.tex} and
    \path{wolf_cor67_maximal_support_restriction.tex}: these describe
    formalization boundaries and their resolution, not a false printed theorem.
\end{itemize}

\section{Conclusion}

At present the formalization-derived record contains five false printed
assertions (Sections~\ref{sec:lorentz},
\ref{sec:pauli-block-truncation}, and the two corrections in
Section~\ref{sec:theorem-6-3}, together with
Section~\ref{sec:unitary-comparison}), one normalization typo
(Section~\ref{sec:ordered-cp-inverse}), one explicit boundary convention
(Section~\ref{sec:radon-nikodym}), and one local correction to a spectrum
shorthand (Section~\ref{sec:two-pure-state-spectrum}). These are the
corrections exposed by the Lean development and its source-faithfulness
audits.

\appendix
\section{Corrected errors in the local LaTeX transcription}
\label{app:transcription}

This appendix is deliberately placed last.  Its entries were found by comparing
the source PDFs with the local transcription; they are \emph{not} claims that a
Lean proof forced a correction to Wolf's statement.  Each was a mathematically
material transcription error and has now been corrected in the local TeX.

\subsection{The Choi marginal: $\tau_B$ versus $\operatorname{tr}_B\tau$}

\textbf{Source context.}  In Chapter~2, immediately after deriving the
coordinate expression for the Choi--Jamiolkowski correspondence in
Eq.~(2.4), Wolf explains the channel--state dictionary.  For a channel
$T:M_d\to M_{d'}$, the first factor is the output factor $A$ and the second is
the untouched input/reference factor $B$.  The PDF (book p.~34) states:
\begin{quote}
``If $T$ is a quantum channel in the Schr\"odinger picture, then $\tau$ is a
density operator with reduced density matrix $\tau_B=\one/d$.''
\end{quote}
Thus $\tau_B$ denotes $\operatorname{tr}_A\tau$.

\textbf{Transcription and correction.}  A previous local transcription at
\path{ch02_representations.tex:113--114} changed the displayed marginal to
$\operatorname{tr}_B\tau=\one/d$.  It has been corrected to retain Wolf's
notation; equivalently, one may spell out
\begin{equation}
  \tau_B=\operatorname{tr}_A\tau=\frac{\one_d}{d}.
\end{equation}

\textbf{Reason.}  Trace preservation is equivalent to the last identity.  The
other marginal satisfies $\operatorname{tr}_B\tau=T(\one)/d$ and is maximally
mixed precisely when $T$ is unital.  Hence the local transcription accidentally
turns a trace-preserving assertion into a stronger and generally false unitality
assertion.

\subsection{The conjugation in the proof of Kramer's theorem}

\textbf{Source context.}  Chapter~3 proves Kramer's theorem for a Hermitian
$H$ commuting with an anti-unitary symmetry $T=\Gamma V$, where $V^T=-V$.
After showing that $V^\dagger\lvert\overline\psi\rangle$ is another eigenvector,
the PDF (book p.~62, Eq.~(3.29)) establishes orthogonality by the chain
\begin{equation}\tag{PDF (3.29)}
 \langle\psi\mid V^\dagger\mid\overline\psi\rangle
 =-\langle\psi\mid\overline V\mid\overline\psi\rangle
 =-\overline{\langle\overline\psi\mid V\mid\psi\rangle}.
\end{equation}
The final conjugation is then identified with
$\langle\psi\mid V^\dagger\mid\overline\psi\rangle$, proving that this scalar
is its own negative.

\textbf{Transcription and correction.}  A previous local TeX version at
\path{ch03_positive_not_completely.tex:515} omitted the outer bar in the final
term.  It has been corrected to the source-faithful and mathematically
necessary term $-\overline{\langle\overline\psi\mid V\mid\psi\rangle}$.

\textbf{Reason.}  Componentwise,
\begin{equation}
 \langle\psi\mid\overline V\mid\overline\psi\rangle
 =\overline{\langle\overline\psi\mid V\mid\psi\rangle}.
\end{equation}
Without the conjugation the equality is false in general, and the final
self-cancellation in the orthogonality argument does not follow.

\subsection{The order of the conjugated factors in Example~5.3}

\textbf{Source context.}  Chapter~5 studies the adjoint of the map
$T(X)=\tfrac12(X^T+\tfrac12\operatorname{tr}(X)\one)$ as an example of a
Schwarz map that need not be $2$-positive.  In the traceless case, the PDF
(book p.~77, Eq.~(5.10)) subtracts the product
\begin{equation}\tag{PDF (5.10), relevant term}
  \frac14\,\overline A A^T
\end{equation}
from $T^*(A^\dagger A)$.

\textbf{Transcription and correction.}  A previous local TeX version at
\path{ch05_schwarz_inequalities.tex:373,377} instead wrote
$A\overline A^{\,T}$.  It has been corrected to $\overline A A^T$, as in the
source.

\textbf{Reason.}  Since $T^*(A)=A^T/2$ for traceless $A$,
\begin{equation}
 T^*(A^\dagger)T^*(A)
 =\frac14(A^\dagger)^T A^T
 =\frac14\,\overline A A^T.
\end{equation}
The transcribed matrix $A\overline A^{\,T}=AA^\dagger$ is generally a
different operator, even though both matrices are positive semidefinite and
the final positivity conclusion in the example survives.

\end{document}
